
We applied projection-only LoRA to Mamba-130m using the standard community configuration --- rank 8, targeting \texttt{in\textbackslash \_proj}, \texttt{out\textbackslash \_proj}, and \texttt{x\textbackslash \_proj}, 1.14\% trainable parameters --- and observed complete failure of classification learning across all four evaluated GLUE tasks. SST-2 accuracy sat at 0.5092 across three identical training epochs (zero-shot: 0.4908); MNLI accuracy degraded from 0.3463 to 0.3234; QNLI oscillated near its 0.5046 baseline; QQP remained at 0.0000. The overall GLUE average improved by less than 0.1pp over zero-shot. The MUST\_WORK gate criterion (SST-2 > 0.70) was not satisfied.

The failure is not a configuration error. LoRA installs correctly, adapters receive gradient updates, and the experiments ran on a functional training pipeline (confirmed by MNLI's monotonic degradation, which demonstrates that the optimizer is updating the model). The failure occurs in the quality of the gradient signal reaching the projection-layer adapters, not in the installation of the adapters themselves.

A critical experimental finding that distinguishes our work from prior characterizations: the experiments used the sequential Python fallback implementation of the Mamba SSM scan, not the custom CUDA kernel. This means the gradient barrier effect does not depend exclusively on CUDA kernel backward-pass properties. The failure manifests under standard PyTorch autograd and likely reflects the mismatch between the gradient signal produced by a randomly-initialized classification head and the gradient landscape conditioned by next-token prediction pretraining through the SSM scan computation.

Our three contributions are: (1) empirical evidence of projection-only LoRA failure on Mamba-130m across all four evaluated GLUE tasks, with precise numerical characterization; (2) verified zero-shot GLUE baselines for \texttt{mamba-130m-hf} reusable by future Mamba adaptation experiments; and (3) a diagnostic signature --- loss oscillation without monotonic decrease, combined with flat or degrading per-epoch accuracy --- for practitioners to detect this failure mode during training.

Three directions follow directly. Exact replication of the MambaPEFT result (90--92\% SST-2) is the highest-priority step, following the ranked ablation order (checkpoint type first). If adaptation is needed at the scan level, \texttt{dt\textbackslash \_proj} LoRA and direct B/C parameter adaptation are the theoretically motivated next steps. And prefix tuning bypasses the scan entirely by adapting at the input embedding level.

Standard LoRA asks Mamba to learn classification through a computation optimized for next-token prediction. The contribution of this work is measuring the failure precisely, describing its signature across multiple tasks, and identifying that the failure manifests even under standard autograd --- pointing to a deeper incompatibility between the SSM pretraining regime and projection-layer classification adaptation.

